\chapter{Background and Related Work}
\label{chap:background_related_work}

\section{Drug--Target Interaction Prediction}

\subsection{Problem Definition}

Drug--target interaction (DTI) prediction aims to determine whether a drug
$d$ interacts with a protein target $p$. Each candidate pair has a binary
label:

\begin{equation}
    y_{dp} \in \{0,1\},
\end{equation}
where $y_{dp}=1$ denotes a known interaction and $y_{dp}=0$ denotes an unknown
or sampled negative interaction. A predictive model estimates:

\begin{equation}
    \hat{y}_{dp}=F(d,p),
\end{equation}
where $\hat{y}_{dp}$ is the predicted probability of interaction.

Computational DTI prediction supports drug repurposing, target identification,
virtual screening, and the investigation of off-target effects. Since
experimentally verified interactions represent only a small fraction of all
possible drug--protein pairs, the task is naturally sparse and imbalanced.

\subsection{Challenges in DTI Prediction}

The first challenge is incomplete supervision. Unobserved pairs are commonly
sampled as negative examples, although some may correspond to interactions
that have not yet been experimentally confirmed. The second challenge is the
heterogeneity of biomedical evidence. Drug structure, protein sequence,
pathways, diseases, side effects, and functional annotations provide
complementary information but differ substantially in representation and
scale.

An effective DTI framework should therefore integrate structural information
with relational biomedical context. It should also model interactions between
drug and protein features rather than treating each feature source
independently.

\subsection{Warm-Start Evaluation Settings}

The current experiments use warm-start evaluation, in which the drugs and
proteins represented in the test pairs are also available in the training
data. The model is evaluated on unseen drug--protein combinations rather than
entirely unseen entities. Two positive-to-negative class ratios are used.

The warm-start 1:1 setting contains one positive interaction for each sampled
negative interaction. This balanced setting provides a controlled assessment
of model discrimination.

The warm-start 1:10 setting contains one positive interaction for every ten
sampled negative interactions. This setting more closely reflects DTI sparsity,
where known interactions form only a small fraction of all possible pairs. Both
ROC-AUC and PR-AUC are reported, with PR-AUC being particularly informative
under the imbalanced 1:10 setting.

\section{Biomedical Knowledge Graph Representation}

\subsection{Biomedical Knowledge Graph Structure}

A biomedical knowledge graph is defined as:

\begin{equation}
    \mathcal{G}
    =
    (\mathcal{E},\mathcal{R},\mathcal{T}),
\end{equation}
where $\mathcal{E}$ is the set of entities, $\mathcal{R}$ is the set of
relation types, and $\mathcal{T}$ is the set of observed triples. Each triple
is represented as:

\begin{equation}
    (h,r,t) \in \mathcal{T},
\end{equation}
where $h,t\in\mathcal{E}$ are the head and tail entities and
$r\in\mathcal{R}$ is the relation type. The direction and type of an edge
express the semantic role of the connected entities.

\subsection{Biomedical Entities and Relations}

Biomedical entities may include drugs, proteins, genes, diseases, pathways,
side effects, biological processes, molecular functions, domains, and protein
families. Relations may represent drug--target interactions, drug--disease
associations, protein--protein interactions, protein--pathway associations,
functional annotations, and domain membership.

This heterogeneous structure provides evidence beyond the observed DTI
matrix. A drug may be characterized by its targets, diseases, side effects,
and therapeutic classes, while a protein may be characterized by pathways,
functions, domains, diseases, and interacting proteins.

\subsection{General Knowledge Graph Embedding Principle}

Knowledge graph embedding (KGE) maps entities and relations into a continuous
vector space. A KGE model defines a scoring function:

\begin{equation}
    f(h,r,t)
\end{equation}
that measures the plausibility of a triple. Training encourages observed
triples to receive higher scores than implausible triples. The learned entity
embeddings provide compact representations of relational context and can be
used for link prediction or supplied to a downstream supervised model.

\subsection{Triple Scoring and Negative Sampling}

For each positive triple $(h,r,t)$, corrupted triples can be generated by
replacing the head or tail entity:

\begin{equation}
    (h,r,t)
    \rightarrow
    (h',r,t),
    \qquad
    (h,r,t)
    \rightarrow
    (h,r,t').
\end{equation}

A margin-based ranking objective separates observed and corrupted triples:

\begin{equation}
    \mathcal{L}_{KGE}
    =
    \sum_{(h,r,t)\in\mathcal{T}^{+}}
    \sum_{(h',r,t')\in\mathcal{T}^{-}}
    \max
    \left(
        0,
        \gamma
        +f(h',r,t')
        -f(h,r,t)
    \right),
    \label{eq:background_kge_loss}
\end{equation}
where $\gamma$ is the margin, $\mathcal{T}^{+}$ contains observed triples,
and $\mathcal{T}^{-}$ contains sampled corrupted triples.

\subsection{Role of KG Embeddings in DTI Prediction}

Within a DTI framework, KGE compresses graph context into drug and protein
representations. For an entity $e$, the KGE stage produces:

\begin{equation}
    \mathbf{z}_{e}
    =
    \operatorname{KGE}(e,\mathcal{G}).
\end{equation}

The embeddings $\mathbf{z}_{d}$ and $\mathbf{z}_{p}$ summarize the relational
contexts of a candidate drug and protein. These embeddings can be combined
with molecular and sequence descriptors in a downstream DTI model. This design
is useful for low-degree entities, whose direct interaction evidence is
limited but whose auxiliary biomedical relations remain informative.

\section{Graph Convolutional Networks}

\subsection{GCN on Undirected Graphs}

For an undirected graph $\mathcal{G}=(\mathcal{V},\mathcal{E},\mathbf{X})$, a
Graph Convolutional Network (GCN) updates node representations by aggregating
normalized neighboring features \cite{kipf2017semi}:

\begin{equation}
    \mathbf{H}^{k+1}
    =
    \sigma
    \left(
        \widehat{\mathbf{A}}\mathbf{H}^{k}\mathbf{W}^{k}
    \right),
    \qquad
    \mathbf{H}^{0}=\mathbf{X},
    \label{eq:background_gcn_update}
\end{equation}
where
$\widehat{\mathbf{A}}=\widetilde{\mathbf{D}}^{-1/2}
(\mathbf{A}+\mathbf{I})\widetilde{\mathbf{D}}^{-1/2}$ is the normalized
adjacency matrix with self-connections, $\mathbf{W}^{k}$ is a learnable
layer-specific transformation, and $\sigma$ is an activation function.
Stacking layers allows each node to incorporate information from increasingly
distant neighborhoods.

\subsection{GCN on Multi-Relational Graphs}

A knowledge graph additionally assigns a relation type and direction to each
edge. A direct multi-relational extension uses a relation-specific
transformation for each neighboring edge:

\begin{equation}
    \mathbf{h}_{v}^{k+1}
    =
    \sigma
    \left(
        \sum_{(u,r)\in\mathcal{N}(v)}
        c_{v,r}^{-1}\mathbf{W}_{r}^{k}\mathbf{h}_{u}^{k}
    \right),
    \label{eq:background_relational_gcn_update}
\end{equation}
where $\mathbf{W}_{r}^{k}$ transforms messages associated with relation $r$
and $c_{v,r}$ is a normalization constant. A separate matrix for every
relation captures edge semantics but causes parameter growth as the relation
vocabulary expands. Directed-GCN adds inverse edges and shares transformations
according to edge direction \cite{MarcheggianiT17}, while R-GCN uses
basis or block-diagonal decomposition \cite{schlichtkrull2018modeling}.
CompGCN instead composes entity and relation embeddings and uses
direction-specific transformations; its update rule is specified in
Chapter~\ref{chap:methodology}.

\section{Existing DTI Prediction Approaches}

\subsection{Descriptor-Based Methods}

Descriptor-based methods formulate DTI prediction as supervised classification
over drug and protein features. Drugs are commonly represented by molecular
fingerprints or chemical descriptors, while proteins are represented by
sequence-derived descriptors. Classical classifiers and neural networks then
learn a decision function from the combined pair representation.

These methods are computationally practical and can generalize through
structural similarity. However, they do not directly exploit heterogeneous
biomedical relations such as disease associations, pathway membership, and
protein--protein interactions.

\subsection{Network and Knowledge-Graph-Based Methods}

Network-based methods use the topology of drug, protein, disease, pathway, and
other biomedical networks. Knowledge-graph-based methods extend this idea by
representing multiple entity and relation types in a unified graph. KGE models
then learn continuous representations that summarize relational compatibility
and graph context.

These approaches provide information beyond direct DTI observations, but the
quality of their representations depends on graph coverage, relation modeling,
negative sampling, and the expressiveness of the selected encoder.

\subsection{Recommendation-Based Methods}

Recommendation-based methods interpret drugs and proteins as two interacting
entity groups, analogous to users and items. Factorization models learn latent
interaction patterns, while neural recommendation models capture nonlinear
relationships among heterogeneous feature groups.

The KGE-NFM framework combines the strengths of knowledge-graph and
recommendation-based approaches \cite{ye2021unified}. KG embeddings provide
relational biomedical representations, structural descriptors provide
intrinsic drug and protein information, and NFM models interactions among
these feature sources. Chapter~\ref{chap:methodology} specifies how this report
implements KGE-NFM with two alternative KG encoders.
